\section{Related work}
\label{sec:related}

\paragraph{Fraud detection on payment data.} Statistical fraud detection has a long literature
across credit cards, telecommunications and money laundering, surveyed by Bolton and Hand
\cite{bolton2002fraud}. Practitioner accounts of card fraud detection stress the problems this
paper meets in synthetic form: extreme class imbalance, drift in fraud behaviour over time, and
metrics that must reflect a limited budget of analyst alerts \cite{dalpozzolo2014lessons}.
Learning under imbalance has its own methods and pitfalls \cite{he2009imbalanced}, and at a base
rate such as this benchmark's \nFraudRateAll{} the receiver-operating curve can look reassuring while precision stays low
\cite{saito2015prc}; we therefore report recall at a fixed false-positive rate and precision beside
its ceiling rather than AUC alone. Our models are standard: gradient-boosted trees in their
XGBoost \cite{chen2016xgboost} and LightGBM \cite{ke2017lightgbm} implementations, and an
Isolation Forest \cite{liu2008iforest} for anomaly scores. We claim no novelty in them.

\paragraph{Mobile money and its fraud.} In Kenya, access to M-PESA has been estimated to have
lifted households out of poverty \cite{suri2016mpesa}. Phipps et al.\
analyse attacks on mobile money systems through ThinSIM add-on cards \cite{phipps2018thinsim}. Lee
et al.\ found that the authentication the prepaid US carriers they studied applied to SIM-swap
requests could be subverted easily, and that the swap alone was enough to compromise accounts on a number of
websites \cite{lee2020simswap}; SIM-swap takeover is one of the benchmark's \nScenarios{} scenarios. Razaq
et al.\ describe SMS- and call-based social-engineering fraud from \nRazaqInterviews{} interviews in Pakistan,
mostly with victims \cite{razaq2021scams}, and Do{\u{g}}an et al.\ survey the scams Kenyan
M-PESA users face and how they detect and respond to them, finding forged M-PESA messages the most
common \cite{dogan2025grift}. With a press account of the ``sent by mistake'' reversal scam
\cite{star2025mpesa}, that work documents the scam the pre-specified variant of
Section~\ref{sec:generalisation} is modelled on.

\paragraph{Synthetic transaction data.} Real transaction data is rarely shareable, which is why
synthetic generators exist. The closest to ours is PaySim, an agent-based simulator of mobile
money transactions calibrated to a sample of a real service's logs, which its paper places in ``an
African country'' without naming it \cite{lopezrojas2016paysim}. Its
log records the client, the action, the recipient, the sum and the change in balances, at an hourly
step; it has no channel, device, geographic or cross-border corridor field. The same authors'
BankSim simulates bank payments between customers and merchants, calibrated with aggregated
transaction data from a bank in Spain \cite{lopezrojas2014banksim}. More recent public benchmarks
take other routes. The Bank Account Fraud suite generates tabular data from an anonymised real
dataset of bank-account-opening fraud, and gives each variant a specific, deliberate bias so that a
method's fairness can be stress-tested beside its performance \cite{jesus2022baf}. Altman et al.\
build an agent-based generator of financial transactions for anti-money laundering, calibrated to
real transactions, and argue that a generator's labels are complete in a way real labels, which miss
undetected laundering, are not \cite{altman2023aml}. BAF's generation from real data is the clearest
contrast with ours, whose fraud behaviour is assumed rather than learned from any real record
(Section~\ref{sec:dataset}). FraudShield-EAC differs in what it
simulates (\nChannels{} channels including USSD and agent banking, devices, coordinates, \nCountries{}
countries and their corridors) and in what it is used for: this paper is less about the generator than about
what its benchmark can and cannot support. \textbf{FraudShield-EAC is not the first public East African fraud dataset.} A structured search of
dataset repositories and the literature, recorded with every query, source and gap
(\texttt{docs/research/c8\_dataset\_search.md} and its second pass,
\texttt{c8\_dataset\_search\_pass2.md}), found earlier ones, among them those below. The Xente Fraud Detection
Challenge on Zindi released real transactions from ``an e-payments, e-commerce, and financial
services company in Uganda'', its training transactions labelled as fraud or not, as a competition
benchmark whose terms
forbid redistribution \cite{zindi2019xente}. MoMTSim's synthetic mobile money transactions, from
Makerere University, are labelled for account takeover (including SIM swaps), refund fraud and fake
credentials, and give Kampala and the Eastern Africa region as the data source location
\cite{azamuke2025momtsim}; produced by a multi-agent simulation of mobile money, it is the closest
neighbour to ours the search found. The search found no public dataset that combines several EAC countries with all
the channels above, but it was bounded: it was run by an AI agent (Section~\ref{sec:disclosure}), and
its records list the sources it could not query and the results it read only in part. That is the
absence of a hit in a bounded search, not evidence that no such dataset exists, and we claim no
novelty on it. We document the dataset with a datasheet \cite{gebru2021datasheets} and the model
with a model card \cite{mitchell2019modelcards}.

\paragraph{Leakage and shortcuts.} Leakage, information about the target that would not be
available at prediction time, is a recurring cause of over-optimistic results in data mining
\cite{kaufman2012leakage} and across machine-learning-based science \cite{kapoor2023leakage}.
Models also learn shortcuts, cues that separate the classes in the data without being the
behaviour of interest \cite{geirhos2020shortcut}. Data with temporal or grouped structure needs
cross-validation that respects it \cite{roberts2017cv,bergmeir2012cv}. Several of the failure
modes in Section~\ref{sec:discussion} are instances of these: rows of one fraud incident leaking
into each other's encoding, and a benchmark whose strongest signal is a single burst feature.

\paragraph{Calibration and explanation.} Scores become probabilities through Platt scaling, a sigmoid fitted to
the classifier's scores \cite{platt2000probabilistic}, or isotonic regression
\cite{zadrozny2002calibration}; Niculescu-Mizil and Caruana compare them \cite{niculescumizil2005calibration}. Calibration error
is usually summarised by binned estimates such as the expected calibration error
\cite{naeini2015bbq,guo2017calibration}, which depend on the binning and can underestimate the
error \cite{kumar2019calibration}; we report the decision-region figure beside the overall one for
that reason. Explanations are Shapley-value attributions \cite{lundberg2017shap}, computed exactly
for tree ensembles by TreeSHAP \cite{lundberg2020trees}.

\paragraph{Evaluation statistics.} We use Hanley and McNeil's standard error for single AUCs
\cite{hanley1982auc}, DeLong's test for correlated AUCs \cite{delong1988auc}, Wilson intervals for
proportions on small counts \cite{wilson1927interval}, and the bootstrap \cite{efron1979bootstrap}
for the gate metrics.

\paragraph{Engineering and research practice.} Machine-learning systems accumulate hidden costs at
their boundaries \cite{sculley2015debt}, and production readiness can be scored by the tests a
system has \cite{breck2017testscore}; the parity suite and evidence stamps of
Sections~\ref{sec:features} and~\ref{sec:repro} are in that spirit. Pre-registration posts the analysis plan to an independent registry before the
research outcomes are observed \cite{nosek2018preregistration}. This work did not do that: its predictions were committed
to its own repository, and Section~\ref{sec:discussion} reports which of them git orders before their
evidence and which it does not.
